\subsection{热力学第一定律. 内能}
\label{sec:app-b-first-law-internal-energy}
\renewcommand{\thestatement}{B.2.\arabic{statement}}
\renewcommand{\theHstatement}{appB.sec02.\arabic{statement}}
\setcounter{statement}{0}

现在考虑一个封闭的均匀热力学系统 (即它不与外界交换物质, 但仍然可以以热或功的形式交换能量). 热力学第一定律告诉我们, 存在一个称为内能并记为 $E$ 的状态变量, 它依赖于温度, 压力和密度, 而且系统在一次变换中的能量变化必须恰好等于从外界提供给系统的能量. 即使系统在宏观层面处于静止状态, 这个内能仍然计入系统中每个分子的动能.

对于封闭系统的一次给定元变换, 有
\[
 \Delta E=\Delta Q+\Delta W,
\]
其中 $\Delta Q$ 是所提供的热量, $\Delta W$ 是变换过程中作用于系统的外力所做的功. 若假定系统除压力外不受其他外力, 则功这一项化为压力力所做的功, 可写成
\[
 \Delta W=-P\Delta V,
\]
其中 $\Delta V$ 是系统在变换过程中发生的体积变化.

\begin{Remark}
\label{rem:app-b-exact-differential}
在上面的能量方程中, $\Delta E$ 表示从初态到终态时状态变量 $E$ 的变化; 重要的是, 这个值不依赖于系统从初态到终态所经过的路径. 然而, $\Delta Q$ 和 $\Delta W$ 的值并不表示状态变量的变化, 并且它们确实依赖于系统在所考虑两个状态之间所沿的路径.

通常把 $\Delta E$ 写成函数 $E$ 的全微分 $dE$ (把 $E$ 看作某一组独立变量的函数), 这是一种标准而有用的记号, 但不能对 $\Delta Q$ 和 $\Delta W$ 作同样处理.

本书中遇到的任何状态函数都使用同样的``微分''记号, 但我们不需要微分形式理论的任何精确内容.
\end{Remark}

\begin{Remark}
\label{rem:app-b-fluid-element-total-energy}
在流体力学框架下, 每个流体微团都含有内能, 因而该流体微团的总能量为
\[
 E_{\mathrm{tot}}=E_{\mathrm{int}}+E_{\mathrm{kin}}+E_{\mathrm{pot}},
\]
其中 $E_{\mathrm{kin}}$ 是系统的宏观动能 (由流动中的平均速度场 $v$ 计算), $E_{\mathrm{pot}}$ 是其势能. 在这个公式中, $E_{\mathrm{int}}$ 是所考虑流体微团中包含的总内能. 实际上, 这个内能项可以表示成内能密度的积分项.
\end{Remark}
